\documentclass[conference,letterpaper,10pt]{IEEEtran}
\usepackage[T1]{fontenc}
\usepackage{cite}
\usepackage{url}
\usepackage[hidelinks]{hyperref}
\hypersetup{pdftitle={Evidence-Bound Recovery for Ambiguous Submissions in Local Generative-AI Pipelines},pdfauthor={}}
\begin{document}
\title{Evidence-Bound Recovery for Ambiguous Submissions in Local Generative-AI Pipelines\\{\large Tool Description}}
\author{\IEEEauthorblockN{Anonymous submission draft}}
\maketitle
\begin{abstract}
A lost submission response can leave a generation client uncertain whether its backend executed a request. We describe evidence-bound recovery tooling for local image-generation pipelines: a versioned evidence exporter, a same-run submission guard, and controlled fault injection with backend observation. The exporter separates reported state, interpretation and oracle evidence. The guard preserves explicitly unknown outcomes without claiming automatic reconciliation. Four fixed Windows/ComfyUI control and response-loss cases were oracle-evaluable. Each control had one submission and one observed execution; each injected case had an unresolved first call and a resolved second call, with two submissions and executions. The second call created a fresh run with a different seed after observed completion. The demonstration establishes bounded execution visibility, not recovery of the original run or fewer executions. Deployment rates and output-quality improvement remain unmeasured.
\end{abstract}
\section{Introduction}

Consider a client that submits an image-generation workflow and loses the
response after the backend accepts it. The backend may continue while the
client lacks the returned job identifier. A second submission creates another
opportunity for execution, but a timeout or a count of HTTP requests cannot
establish what the backend actually did. The relevant question is which
observations justify the next recovery action and the subsequent report of
that action.

Uncertain outcomes are an established distributed-systems problem. Remote
procedure call designs distinguish call semantics under communication failure
\cite{birrell1984}, while idempotent API designs use caller-provided identifiers and explicit
service behavior to make retries safe \cite{aws}. A local workstation does not remove
the distinction between request transmission and backend execution. Our tool
addresses how that distinction is represented and inspected in a particular
local generation pipeline; it does not introduce a new exactly-once protocol.

\texttt{local-gpu-imagegen} combines persistent run records, generation orchestration,
a ComfyUI adapter and artifact handling. We describe tooling that preserves the
client's evidence, applies versioned interpretation rules, and records backend
observations independently of the client's returned status. Its product guard
is deliberately narrower: an explicitly unknown submission blocks further
submission through the same unresolved run at the reviewed entry points.

The tool description has three parts: an evidence representation and read-only
exporter; a same-run guard whose unresolved state remains visible; and CPU and
Windows fault-observation workflows with separate evidence boundaries. The
Windows demonstration tests visibility across the client--backend boundary.
It does not establish that the guard reduces execution counts in that protocol.
The contribution is an inspectable integration of evidence interpretation,
run-scoped control and backend observation. Section 6 compares its scope with
idempotent APIs, workflow execution and tracing. An exhaustive nearest-tool
review remains open; no first-of-its-kind claim is made.

\section{Failure model and evidence semantics}

A \textbf{product call} is an invocation by the research client. A \textbf{run} is the
persistent product object created by \texttt{local\_gpu\_start\_run}. A \textbf{submission} in
the Windows results is a proxy-observed \texttt{POST /prompt}; a call that fails before
that boundary need not produce a submission. An \textbf{execution instance} is an
observation bound to backend execution events. CPU fixtures instead observe
the synthetic worker entry. These units are not interchangeable.

F00 is a normal control. F02 suppresses the first accepted submission response,
creating a window in which the client lacks an acknowledged job identifier.
F03, covered by the CPU fixtures rather than the Windows campaign reported
here, concerns loss of completion information after a job identifier is known.
An explicit rejection and a missing response have different evidential meanings.
Neither should be silently converted into a backend execution outcome.

The evidence model distinguishes reported state, interpreted state and oracle
state. In addition, the Windows research client emits its own call-level labels.
Its \texttt{resolved} label is derived from a generated product result after checking
for artifact hashes. Its \texttt{unresolved} label can arise from several client errors.
Thus a research-client label alone does not prove that a durable product
manifest has a particular recovery state. In F02, the proxy's accepted-response
loss and the oracle binding provide the additional context.

\section{Tool design}

\subsection{Evidence-preserving export}

The exporter accepts a record or run manifest and produces a separate normalized
record. It copies reported state, records the mapping version and reasons, and
retains supplied oracle state separately. The schema includes source identity,
request and input digests, backend identity, job identifier, artifact hash and
validator version. Missing fields carry explicit reasons instead of fabricated
values. The inspected mapping is \texttt{research-normalization-v2}.

Interpretation checks more than an execution-success label. Job and artifact
bindings, approval and validation evidence, and recovery state constrain the
\texttt{execution\_verified} decision. An unresolved recovery state prevents this flag
from becoming true even when other evidence reports successful execution.
This permits the record to express backend completion without asserting that
the product's recovery obligation has been discharged. The exporter does not
invoke the engine, retry a request or update the source run. Its output writer
uses exclusive file creation. These are implementation properties, not a
claim of complete filesystem security.

\subsection{Recovery guard and its scope}

\texttt{AssetRunEngine} and \texttt{RunStore} preserve the explicit transport marker
\texttt{submission\_outcome=unknown} as an unresolved attempt. The reviewed submission
entry checks that unresolved attempt before permitting further generation.
The guard applies within the same run, including attempts with a different
idempotency key. The corresponding inspection action is \texttt{get\_run}.

This guard does not discover a missing backend job identifier or reconcile an
unknown job automatically. A newly created run lies outside the prior run's
guard. Existing recovery for a known job in a two-stage operation is a separate
path and is not attributed to this no-job-ID guard. The resulting tradeoff is
explicit: the guard withholds another submission through that run while the
required evidence is unavailable; it does not guarantee eventual completion.


\subsection{Controlled observation on Windows}

Figure~\ref{fig:architecture} shows the product and research boundaries.
The research harness adds a transport shim and loopback fault proxy. The shim
routes prompt transport through the proxy while retaining the actual backend
endpoint in the product's identity client. F02 drops the first response only
after backend acceptance. This arrangement is a research instrument, not a
native ComfyUI recovery feature.

The observer connects before the product invocation and uses the operation key
as its ComfyUI client identifier, matching the product's event-routing identity.
\texttt{ComfyUIEventOracle} requires an execution-start event, a terminal event and
completed history bound to a prompt identifier. It derives an execution-instance
identifier using the backend boot identity, prompt identifier and observed
start sequence. Its independence is from the product return value: it still
trusts the backend's event and history interfaces, and it is not a physical GPU
monitor or a crash-durable audit log.

The controller permits the second F02 product call only after the first
execution has a bindable completion. The client creates a new run on each call
and uses \texttt{seed = 4100 + call\_index}. This is observation-gated fresh-run
generation, not automatic reconciliation of the original run. Matching frozen
request digests does not establish full prompt equality: the inspected digest
function excludes the seed and does not use its plan argument.



\begin{figure}[t]
\centering
\setlength{\fboxsep}{5pt}
\begin{tabular}{c}
\fbox{Product engine, run store and adapter}\\[3pt]
$\downarrow$ prompt transport\\[3pt]
\fbox{Research shim and fault proxy}\\[3pt]
$\downarrow$ accepted submission\\[3pt]
\fbox{ComfyUI backend}\\[3pt]
$\downarrow$ events and history\\[3pt]
\fbox{Observer and second-call gate}
\end{tabular}
\caption{Product and research boundaries. The controller's second call creates
 a fresh run. The exporter separately normalizes reported records; it does not
 trigger the controller or change the source run.}
\label{fig:architecture}
\end{figure}

\section{Evidence and demonstration}

\subsection{CPU components and deterministic fixtures}

The CPU harness separates submission events from synthetic worker-entry events.
Its fixtures cover normal operation and response loss before and after job-ID
acquisition, including single-stage and two-stage paths. Oracle self-checks
exercise cases in which submission and execution counts differ and in which a
shared prompt identifier does not imply a single execution. These are fixed
semantic checks, not randomly sampled deployments.

Historical command output records 40 research component tests passing in
7.110 seconds using Python unittest discovery, with \texttt{OK} and exit code 0. The preparation report identifies the runtime
as Windows Python 3.15.0a8. A later macOS check targeted the transport shim
and campaign controller: seven tests passed in 0.003 seconds, followed by
successful compilation and whitespace checks. These are distinct stages;
the earlier 40-test result is not a full-suite result for the final shim.
Neither count is an experimental sample size. Exact immutable source binding
of the historical test worktrees remains incomplete.

Earlier paired CPU results are documented separately. They must not be
pooled with the Windows cases: the synthetic same-run path and the fresh-run
Windows client do not exercise the same recovery scope.

\subsection{Four controlled Windows cases}

We denote the frozen baseline product revision by B2 and the revision with
the same-run ambiguity guard by W3. The endpoint-preserving campaign report
records both revisions under F00 and
F02. All four fixed cases were oracle-evaluable; the two F00 cases were controls
and the two F02 cases had confirmed response-loss injection. Table~\ref{tab:cases} reports
the observations without converting these counts into a success probability.


\begin{table}[t]
\caption{Four fixed cases, checked against retained report fields. S/E are
proxy submissions/bound execution instances; U/R are unresolved/resolved
research-client calls. All four cases were oracle-evaluable.}
\label{tab:cases}
\centering
\begin{tabular}{llccc}
\hline
System & Case & Calls & S/E & Call sequence\\
\hline
B2 & F00 & 1 & 1/1 & R\\
W3 & F00 & 1 & 1/1 & R\\
B2 & F02 & 2 & 2/2 & U $\to$ R\\
W3 & F02 & 2 & 2/2 & U $\to$ R\\
\hline
\end{tabular}
\end{table}

The campaign totals were six calls, six proxy submissions and six bound
execution instances. The F02 second call used a fresh run and changed seed.
Counts were read from the retained campaign report; raw events have not been
independently reconstructed in this manuscript revision. No rate is estimated.

The report identifies an RTX 5070 Ti Laptop GPU and ComfyUI v0.30.0 at a pinned source revision, with model and workflow hashes
checked against the campaign configuration. These are recorded environment
facts for this campaign. They do not establish behavior across other models,
backend versions or machines.

Both systems produced two observed executions in F02. This demonstrates that
the research observer could retain execution evidence despite the first call's
ambiguous outcome. It does not demonstrate an execution-count advantage for W3.
Because the second call changed run and seed, the table also does not establish
a duplicate execution of an identical generation request.

\subsection{Resolution labels require an explicit aggregation rule}

A resolved second call does not erase the first call's unresolved result or
show that its original run was reconciled. Read-only inspection of the retained
campaign report confirmed that both F02 records carry \texttt{resolved=1} and
\texttt{unresolved=1}. These are independent any-call flags: across the four cases,
\texttt{resolved} sums to four and \texttt{unresolved} to two. They are overlapping counts,
not mutually exclusive outcomes. The earlier Markdown summary's zero unresolved
cases is therefore not the controller aggregation. Table~\ref{tab:cases} preserves the call
sequence; the evidence addendum records this correction without rewriting the
historical report. No recovery-success rate is inferred.

There is also a counting distinction within the observer export. Successive
oracle decisions retain cumulative raw-event snapshots, so summing their event
counts can count earlier events again. The controller's execution count uses
distinct bound execution-instance identifiers. A raw-evidence audit must check
those bindings, not substitute aggregate event counts.

\section{Reproducibility and limitations}

The source and retained reports support inspection of the protocol. Frozen B2
and W3 product revisions are distinct from the research harness revision;
a reproduction must preserve all three identities. Historical commands and
source identifiers are recorded in the companion evidence ledger. A fresh
clean-checkout reproduction has not been performed for this revision.

Raw Windows campaign JSONL, private configurations, model files and generated
images remain Windows-local according to the retained report. Read-only follow-up
inspected the retained campaign-report fields and computed its file hash.
This did not independently reconstruct raw WebSocket events or verify image
content. A sanitized
per-call event/history manifest and its binding checks are still required for
independent inspection of Table~\ref{tab:cases}. The availability of source and a summary
is not equivalent to a complete distributable reproduction package.


Artifact evidence has its own boundary. The research client computes content
hashes for returned artifacts, whereas oracle-exported path fingerprints hash
metadata strings. A path fingerprint is not a content hash. Completed backend
history and a content hash do not establish visual quality, task acceptance or
finalization. Earlier protected-pixel or mask-validation results from a
different pilot are not imported into this campaign.

Observation depends on receiving and binding backend events within the harness
window. Missing events limit evaluability; they do not establish non-execution.
The same-process CPU observer and backend-origin Windows observer have different
failure boundaries, neither of which establishes crash-complete auditability.
The fixed case set contains no basis for deployment failure frequency, natural
duplicate-execution frequency, performance improvement or broad model
robustness. No GPU run, model download or new experiment was performed for this
manuscript revision.

\section{Related work and discussion}

RPC failure semantics provide the conceptual starting point: a failed
communication does not necessarily settle whether a remote operation ran \cite{birrell1984}.
Idempotent API design addresses safe retries through an explicit contract,
including caller-provided request identifiers \cite{aws}. Our tool relies on neither
a prompt identifier nor a client idempotency key as sufficient evidence of
backend deduplication. Its role is to preserve and qualify observations when
such a guarantee has not been established for the evaluated path.

The practical distinction is between observation and control. Backend events
can establish that work completed while the client's original run remains
unresolved. Conversely, a conservative same-run guard can block a further
submission without resolving that uncertainty. The Windows harness crosses a
different boundary by creating a new run after observation. Keeping these
scopes separate explains why the demonstration can be oracle-evaluable without
showing automatic recovery or a W3 reduction in executions.

Temporal's Activity documentation describes the same lost-completion boundary:
a worker may finish an Activity and crash before reporting completion, leading
to a retry. It recommends idempotent Activities \cite{temporal}. This is an important
comparison because durable workflow state alone does not eliminate uncertainty
about external side effects. Our same-run guard withholds a submission under
explicit uncertainty, but supplies no equivalent general workflow-execution or
automatic reconciliation guarantee.

OpenTelemetry represents operations through spans and uses trace context to
relate observations across a request \cite{otel}. Such correlation provides a useful
representation for a distributed execution path. In this tool, the additional
application-specific rule is which observations permit an execution claim:
backend start and terminal events must bind to completed history, and unresolved
product recovery can still veto \texttt{execution\_verified}. Correlation identifiers
alone are not used as execution counts.

Toxiproxy provides configurable TCP connection fault injection for testing \cite{toxiproxy}.
Our research proxy uses an application-level acceptance condition to select
the response to suppress, while the observer separately checks execution. This
comparison does not show that Toxiproxy could not support an equivalent setup;
no comparative implementation or usability measurement was performed.

This is a documentation-based comparison of responsibilities, not a performance
benchmark or a feature-absence survey. The practical contribution is making the
relationship between client ambiguity, backend observation and permitted claims
inspectable for a local generation workflow. A broader academic and
reconciliation-system nearest-work comparison remains incomplete.
 No superiority or exhaustive novelty claim follows from the
selected sources.

\section{Conclusion}

The tool preserves ambiguous client outcomes alongside separately bound
execution evidence and implements a conservative guard within an unresolved
run. The controlled Windows cases demonstrate execution visibility under lost
acceptance responses, while exposing the distinction between a new successful
call and reconciliation of the original run. These bounded observations support
an evidence-preserving tool description. They do not establish automatic
unknown-job recovery, deployment-wide reliability or natural duplicate-execution
rates.

\bibliographystyle{IEEEtran}
\bibliography{references}
\end{document}
